\chapter{Commodity Spreads}\label{ch:s2:commodity-spreads}

A refinery earns the difference between the crude it buys and the products it sells. A trader who holds that difference, rather than oil, is betting on
refining margins, storage and the seasons, not on the price of oil. From 2006 to 2026 the New York Harbor 3-2-1 crack spread, computed from EIA spot
prices, averaged \$21.30 a barrel and reverted with a half-life of about 51 trading days. A rule fading its deviations earned a Sharpe ratio of 0.48
with a correlation of $-0.02$ to crude's daily moves. On this chapter's synthetic market, where the crack has a planted season and a planted
mean reversion, the same rule earns a Sharpe ratio of 1.05, with a correlation of $-0.12$ to crude. The build is \texttt{firm.commspread}.

\section{Calendar spreads again}

A calendar spread (Book~1, chapter~19) is the simplest commodity spread: one commodity at two dates, priced by storage and convenience yield. This
chapter's spreads join different commodities, grades or places. Each is priced by a physical activity that converts one into the other: refining,
generating, crushing, shipping, blending. The activity puts bounds on the spread. When the crack is too narrow, refiners cut runs and product
inventories fall; when too wide, they run harder. That is the economic reason for the mean reversion a trader fades, and the reason it can fail when
the activity itself is blocked.

\begin{definition}[Processing-spread trade]\label{def:s2:commodity-spreads:processing}
A \emph{processing-spread trade}\index{processing-spread trade} is a position in an input commodity against the outputs made from it, in the ratio the
process yields (crude against gasoline and diesel, gas against power at a heat rate, soybeans against meal and oil), that profits from changes in the
processing margin rather than in the commodities' price level.
\end{definition}

Girma and Paulson found crude oil, gasoline and heating oil futures prices cointegrated, with stationary spreads between crude and its products. They
found historically profitable and statistically significant risk arbitrage in three popular petroleum spreads, though they could not be certain that
these opportunities still existed. \texttt{firm.commspread} builds the spreads by ratio
(\cref{lst:s2:commodity-spreads:build}); with gasoline at \$2.50 a gallon, diesel at \$3.00 and crude at \$80 a barrel, the 3-2-1 crack is \$32.

\section{Quality and location}

\begin{definition}[Location spread]\label{def:s2:commodity-spreads:location}
A \emph{location spread}\index{location spread} is the price difference between the same or nearly the same commodity at two delivery points, bounded
in normal times by the cost of moving it between them and widened without bound when the transport capacity is full.
\end{definition}

\begin{definition}[Quality spread]\label{def:s2:commodity-spreads:quality}
A \emph{quality spread}\index{quality spread} is the price difference between two grades of one commodity (light sweet against heavy sour crude,
high- against low-protein wheat), set by the cost of processing the lower grade into the products of the higher one.
\end{definition}

Brent is a seaborne crude priced in the North Sea; WTI is delivered inland at Cushing, Oklahoma. From 2006 to 2026 Brent traded on average \$4.76
above WTI, with a standard deviation of \$5.97 (\cref{fig:s2:commodity-spreads:real}). The EIA attributed the record spread of 2011, when Brent
averaged about \$111 and WTI about \$95, to transportation bottlenecks near Cushing. In 2017 the EIA wrote that near-term changes in the spread
generally come from changes in pipeline capacity or in US crude production. A \omterm{def:s2:commodity-spreads:location}{location spread} is a view on pipes.

The daily data hold two warnings. On 20 April 2020 WTI's spot price was $-\$36.98$, and Brent stood \$54.34 above it, the series' maximum: a delivery
point that could not take more crude priced it below zero. And the Brent-WTI spread's daily changes have an autocorrelation of $-0.36$. Brent is
assessed at the London close and WTI at New York's, so each day's spread mixes two times, and part of its change is undone the next day. A rule
fading the spread's deviations from its trailing mean earned a Sharpe ratio of 0.85 without April 2020. Entered a day later, it earned 0.32.

\begin{omfigure}
\begin{tikzpicture}
\begin{axis}[width=11cm,height=5.4cm,xlabel={\small year},ylabel={\small \$/bbl, monthly mean},tick label style={font=\footnotesize},
  xmin=2006,xmax=2027,ymin=-10,ymax=70,grid=major,grid style={gray!25},table/col sep=comma,
  x tick label style={/pgf/number format/1000 sep={}},legend style={font=\scriptsize,at={(0.5,0.97)},anchor=north}]
\addplot[omThm,thick] table[x=year,y=crack]{figdata/strategies-2/21-commodity-spreads/real.csv};
\addplot[omDef,thick] table[x=year,y=brentwti]{figdata/strategies-2/21-commodity-spreads/real.csv};
\legend{New York Harbor 3-2-1 crack,Brent minus WTI}
\end{axis}
\end{tikzpicture}

\omcaption{Monthly means of the New York Harbor 3-2-1 crack spread and of Brent minus WTI, from EIA daily spot prices, June 2006 to September 2026.
Data: EIA via FRED; \texttt{s2\_fetch\_spreads.py}.}\label{fig:s2:commodity-spreads:real}
\end{omfigure}

\section{Processing spreads: crack, spark, crush}

The synthetic crack in \texttt{firm.commspread} has four parts: a mean of \$20, a seasonal cycle of \$5 amplitude that peaks in late April, before the
summer driving season, a deviation that reverts with a half-life of 40 days, and a planted lag: when crude jumps, products follow over weeks, so the
crack first narrows. Its daily changes have a correlation of $-0.36$ with crude's, close to the $-0.38$ of the real 3-2-1 crack. Fading the crack
beyond 1.5 standard deviations of its trailing year (\cref{lst:s2:commodity-spreads:fade}):

\begin{center}
\small
\begin{tabular}{@{}lccc@{}}
\toprule
nine years, \$0.10/bbl per unit traded & \$/bbl a year & Sharpe ratio & correlation with crude\\
\midrule
fade the crack & 15.4 & 1.05 & $-0.12$\\
fade the crack, season removed & 6.2 & 0.44 & $-0.07$\\
long the crack February--April & 7.4 & 0.78 & $-0.15$\\
fade the \omterm{def:s2:commodity-spreads:location}{location spread} & 3.8 & 0.48 & 0.01\\
fade the \omterm{def:s2:commodity-spreads:quality}{quality spread} & 2.6 & 0.98 & 0.01\\
equal-risk combination & & 1.69 & $-0.13$\\
\bottomrule
\end{tabular}
\end{center}

Removing each month's average before fading cuts the crack trade by more than half: much of what the plain rule fades is the season. The two crack
rules are one bet on the same seasonal cycle. On the real gasoline crack, the change from the start of February to the end of April was positive in 14
of the 20 years from 2007 to 2026, averaging \$6.17 with a $t$-statistic of 2.8; the worst year lost \$11.97.

The spark spread is power less gas at a heat rate: at \$50 per megawatt-hour, gas at \$3 per million British thermal units and a heat rate of 7, it
is \$29. A trader with a view of which plants will set the price, and at what heat rate, trades the spread at that rate. The crush is soybean meal
and oil less the beans they come from. Simon found that deviations of the crush from its long-run equilibrium were transitory from 1985 to 1995, with
strong seasonality and a tendency to revert toward its recent five-day average, and that trading rules based on this would have been profitable.

\section{Risk in spread books}

A spread has lower volatility than its legs, and a spread book looks well diversified: the combination above has a Sharpe ratio of 1.69. Three risks
hide in it. First, the spreads' reversion depends on the physical activity that links the legs, and when that activity is blocked the spread trends.
The synthetic \omterm{def:s2:commodity-spreads:location}{location spread} earned \$35.3 a barrel over four years by fading, then a planted pipeline bottleneck widened it by \$15 over six months.
The fade lost \$14.1 in those months and \$9.3 over the bottleneck's life, before earning \$8.4 afterwards
(\cref{fig:s2:commodity-spreads:books}).

Second, a spread is two or more futures positions, each margined, each with its own liquidity, delivery and expiry; the thinner leg sets the cost.
Third, the ratio in the trade is a convention, not the process: a refinery's yields differ from 3-2-1, and the crush's 44 and 11 pounds are averages.
A book that fades several spreads also shares their common drivers. On EIA spot data, the crack fade and the Brent-WTI fade both had their best day
of the sample on 20 or 21 April 2020, when WTI's price went negative; in such a week a book of spreads is one position.

\begin{omfigure}
\begin{tikzpicture}
\begin{axis}[width=11cm,height=5.4cm,xlabel={\small year},ylabel={\small cumulative \$/bbl per unit},tick label style={font=\footnotesize},
  xmin=1,xmax=10,ymin=-10,ymax=150,grid=major,grid style={gray!25},table/col sep=comma,
  legend style={font=\scriptsize,at={(0.03,0.97)},anchor=north west}]
\fill[gray!15] (axis cs:5,-10) rectangle (axis cs:6.17,150);
\addplot[omThm,thick] table[x=year,y=crack]{figdata/strategies-2/21-commodity-spreads/books.csv};
\addplot[gray,thick,dashed] table[x=year,y=seasonal]{figdata/strategies-2/21-commodity-spreads/books.csv};
\addplot[omDef,thick] table[x=year,y=location]{figdata/strategies-2/21-commodity-spreads/books.csv};
\addplot[black,thick] table[x=year,y=quality]{figdata/strategies-2/21-commodity-spreads/books.csv};
\legend{fade the crack,February--April crack,fade the location spread,fade the quality spread}
\end{axis}
\end{tikzpicture}

\omcaption{Cumulative P\&L of four synthetic spread rules, in dollars per barrel per unit of spread; the shaded band is the planted pipeline bottleneck
that widens the \omterm{def:s2:commodity-spreads:location}{location spread}. Data: \texttt{s2\_commspread.market}.}\label{fig:s2:commodity-spreads:books}
\end{omfigure}

\section{Strategy files}

\begin{strategyfile}{Crack-spread mean reversion}\label{strat:s2:commodity-spreads:crack}
\sfield{Who pays you, and why} Refiners and product buyers hedging at the extremes; refiners' run changes close the gap with a lag.
\sfield{Instruments and venues} Crude, gasoline and diesel futures in the 3-2-1 ratio; exchange-listed crack spread contracts.
\sfield{Signal} The crack's deviation from its trailing mean, with or without its season.
\sfield{Sizing and execution} Traded as one spread order where the exchange allows; sized by the spread's volatility.
\sfield{Costs} Three legs; the thinnest sets the cost.
\sfield{How it dies} Refinery outages and structural shifts in product demand that move the mean.
\sfield{Horizon, capacity, infrastructure} Weeks to months.
\sfield{Backtest honestly} Futures, not spot; synchronous closes; the season separated from the reversion.
\sfield{Sources} Girma and Paulson (1999); EIA spot: Sharpe ratio 0.48 (0.35 without April 2020); this chapter: 1.05 on a planted crack.
\end{strategyfile}

\begin{strategyfile}{Spark spread with a heat-rate view}\label{strat:s2:commodity-spreads:spark}
\sfield{Who pays you, and why} Generators and retailers hedging; the market's view of which plants set the price.
\sfield{Instruments and venues} Power and gas futures or forwards for the same period and region.
\sfield{Signal} The trader's heat rate for the marginal plant against the market-implied one.
\sfield{Sizing and execution} Gas at the heat rate per megawatt-hour of power.
\sfield{Costs} Power's wide bid-ask; shape and location mismatch.
\sfield{How it dies} New capacity, outages, weather.
\sfield{Horizon, capacity, infrastructure} Months to seasons; a fundamental power model.
\sfield{Backtest honestly} The plant stack as it was at the time.
\sfield{Sources} No performance figure verified.
\end{strategyfile}

\begin{strategyfile}{Soybean crush trade}\label{strat:s2:commodity-spreads:crush}
\sfield{Who pays you, and why} Processors' margin hedging; crushing capacity closes wide margins.
\sfield{Instruments and venues} Soybean, meal and oil futures (the board crush).
\sfield{Signal} Deviations from a seasonal equilibrium; short-term reversion.
\sfield{Sizing and execution} 44 pounds of meal and 11 of oil per bushel, in contract units.
\sfield{Costs} Three legs.
\sfield{How it dies} Changes in demand for oil (biofuels) or meal that move the equilibrium.
\sfield{Horizon, capacity, infrastructure} Days to months.
\sfield{Backtest honestly} The trend in the equilibrium; out-of-sample periods.
\sfield{Sources} Simon (1999).
\end{strategyfile}

\begin{strategyfile}{Brent-WTI location spread}\label{strat:s2:commodity-spreads:location}
\sfield{Who pays you, and why} Shippers and refiners paying for access to the other market; the spread pays for pipes and tankers.
\sfield{Instruments and venues} Brent and WTI futures.
\sfield{Signal} The spread against transport costs and capacity.
\sfield{Sizing and execution} One contract against one.
\sfield{Costs} Low; roll mismatch between the two contracts' expiries.
\sfield{How it dies} A bottleneck that widens the spread for a year; delivery-point squeezes.
\sfield{Horizon, capacity, infrastructure} Weeks to years; pipeline and export data.
\sfield{Backtest honestly} Synchronous prices: spot assessments at different closes create false reversion.
\sfield{Sources} EIA (2012, 2017); EIA spot: Sharpe ratio 0.85, 0.32 a day late.
\end{strategyfile}

\begin{strategyfile}{Sweet-sour quality spread}\label{strat:s2:commodity-spreads:quality}
\sfield{Who pays you, and why} Refiners who can or cannot process sour crude; complex refineries pay less for it.
\sfield{Instruments and venues} Grade differentials, swaps or physical cargoes.
\sfield{Signal} The differential against the value of the extra processing.
\sfield{Sizing and execution} Barrel for barrel.
\sfield{Costs} Illiquid differentials.
\sfield{How it dies} Changes in supply of one grade (OPEC cuts of sour crude); refinery conversions.
\sfield{Horizon, capacity, infrastructure} Months.
\sfield{Backtest honestly} Assessed prices, not traded ones.
\sfield{Sources} No performance figure verified; this chapter: 0.98 on a planted spread.
\end{strategyfile}

\begin{strategyfile}{Seasonal gasoline crack}\label{strat:s2:commodity-spreads:season}
\sfield{Who pays you, and why} The switch to summer-grade gasoline and the driving season; refiners' maintenance in spring.
\sfield{Instruments and venues} Gasoline against crude futures, summer contracts.
\sfield{Signal} The calendar.
\sfield{Sizing and execution} Long from February to April.
\sfield{Costs} Two legs, rolled.
\sfield{How it dies} Well-known seasonality priced into the forward curve; a demand shock.
\sfield{Horizon, capacity, infrastructure} Three months a year.
\sfield{Backtest honestly} Futures for the summer month, whose price already includes the season, not spot.
\sfield{Sources} EIA spot: up in 14 of 20 years, mean \$6.17, $t$ 2.8.
\end{strategyfile}

\section{Tutorial: not the price of oil}\label{tut:s2:commodity-spreads:tut}

\begin{tutorial}
\textbf{Goal.} Build spreads by ratio, simulate a crack with a season and a lag and a \omterm{def:s2:commodity-spreads:location}{location spread} with a bottleneck, fade them, and compare with
the EIA spot record. \textbf{End state:} the table and two figures.
\begin{tutsteps}
\item \textbf{Spreads by ratio}.
\omcode{code/firm/commspread/firm_commspread.py}{31}{41}{The 3-2-1 crack, the spark spread and the board crush.}\label{lst:s2:commodity-spreads:build}
\item \textbf{The fade and the season}.
\omcode{code/firm/commspread/firm_commspread.py}{107}{141}{Fading the deviation from the trailing mean, and the seasonal rule.}\label{lst:s2:commodity-spreads:fade}
\item \textbf{Run} \texttt{table()}, \texttt{bottleneck()}, \texttt{bands()}, \texttt{fig\_commspread.py} and, once,
\texttt{s2\_fetch\_spreads.py}.
\end{tutsteps}
\textbf{What to change next.} Add asynchronous closing times to the synthetic \omterm{def:s2:commodity-spreads:location}{location spread} and measure the false reversion; estimate the season
from past years only; stop the location fade when transport capacity is full.
\end{tutorial}

\section{Build: commodity spreads}\label{bld:s2:commodity-spreads:commspread}

\begin{build}
\bfield{Purpose} Spread construction by ratio; a synthetic market of crude, crack, location and \omterm{def:s2:commodity-spreads:quality}{quality spreads}; fade and seasonal rules.
\bfield{Interface} \texttt{crack\_321}, \texttt{spark}, \texttt{board\_crush}, \texttt{SpreadConfig(\dots)}, \texttt{simulate\_spreads(cfg)},
\texttt{fade(x, cfg, month, delay)}, \texttt{seasonal(x, month, cfg)}.
\bfield{Rules} P\&L per unit of spread from the next day's change; costs per unit traded; trailing windows only.
\bfield{Acceptance tests} \texttt{code/firm/commspread/tests/}: the three spreads by hand; the fade's positions and P\&L by hand; the seasonal
window; the bottleneck's shape and the crack's negative response to crude.
\bfield{Stretch} Futures curves for each leg; a refinery model with run cuts.
\end{build}

\begin{omsources}
\item P.~B.~Girma and A.~S.~Paulson, ``Risk arbitrage opportunities in petroleum futures spreads'', \emph{Journal of Futures Markets} 19(8), 1999.
\item D.~P.~Simon, ``The soybean crush spread: empirical evidence and trading strategies'', \emph{Journal of Futures Markets} 19(3), 1999.
\item US Energy Information Administration, ``2011 brief: energy commodity price trends varied widely during 2011'', \emph{Today in Energy}, 2012.
\item US Energy Information Administration, ``Transportation constraints and export costs widen the Brent-WTI crude oil price spread'',
\emph{Today in Energy}, 15 November 2017.
\item US Energy Information Administration, daily spot prices of WTI, Brent, New York Harbor gasoline and diesel, via FRED, accessed 25 September
2026.
\end{omsources}

\section{Exercises}

\begin{exercise}[$\star$]\label{exo:s2:commodity-spreads:1}
Gasoline is \$2.50 a gallon, diesel \$3.00 and crude \$80 a barrel. What is the 3-2-1 crack?
\end{exercise}

\begin{exercise}[$\star$]\label{exo:s2:commodity-spreads:2}
Power is \$50 per megawatt-hour and gas \$3 per million British thermal units. What is the spark spread at a heat rate of 7, and at 7.5?
\end{exercise}

\begin{exercise}[$\star$]\label{exo:s2:commodity-spreads:3}
Soybeans are \$10 a bushel, meal \$300 a short ton and oil 50 cents a pound. What is the board crush?
\end{exercise}

\begin{exercise}[$\star\star$]\label{exo:s2:commodity-spreads:4}
Why do the daily changes of the crack spread correlate negatively with crude's?
\end{exercise}

\begin{exercise}[$\star\star$]\label{exo:s2:commodity-spreads:5}
Why does the Brent-WTI fade lose most of its Sharpe ratio when entered a day later?
\end{exercise}

\begin{exercise}[$\star\star$]\label{exo:s2:commodity-spreads:6}
What economic force makes the crush spread's deviations transitory?
\end{exercise}

\begin{exercise}[$\star\star\star$]\label{exo:s2:commodity-spreads:7}
\emph{Coding.} Run \texttt{bands()}. How do the crack fade's Sharpe ratio and time in the market change with entry bands of 1.0, 1.5 and 2.0
standard deviations, and what would choosing the best one after the fact do to the backtest?
\end{exercise}

\begin{exercise}[$\star\star\star$]\label{exo:s2:commodity-spreads:8}
\emph{Find the flaw.} ``Our location fade made \$35 a barrel in four years; the spread always comes back to transport cost, so we doubled it.''
\end{exercise}

\section{Problem: Not the Price of Oil}

\begin{problem}[{Weekend problem --- commodity spreads}]\label{pb:s2:commodity-spreads:1}
The chapter's synthetic spreads and the EIA record.

\medskip\noindent\textbf{Part I --- Spreads.}
\begin{enumerate}
\item Define a \omterm{def:s2:commodity-spreads:processing}{processing-spread trade}, a \omterm{def:s2:commodity-spreads:location}{location spread} and a \omterm{def:s2:commodity-spreads:quality}{quality spread}.
\item What bounds each spread, physically?
\item What did Girma and Paulson find?
\item What did Simon find for the crush?
\end{enumerate}

\medskip\noindent\textbf{Part II --- The record.}
\begin{enumerate}[resume]
\item Give the real 3-2-1 crack's mean, half-life and correlation with crude.
\item What happened to Brent-WTI in 2011 and in April 2020?
\item Why do the daily Brent-WTI changes revert, and what does it do to a backtest?
\item What did the real gasoline crack do from February to April?
\end{enumerate}

\medskip\noindent\textbf{Part III --- The synthetic market.}
\begin{enumerate}[resume]
\item Describe the synthetic crack.
\item Give the table of rules.
\item Why does removing the season halve the crack fade?
\item What did the bottleneck do to the location fade?
\end{enumerate}

\medskip\noindent\textbf{Part IV --- The verdict.}
\begin{enumerate}[resume]
\item State the \emph{named result}: the crack-spread trade's Sharpe ratio and its correlation with crude.
\item Why is the combination's Sharpe ratio of 1.69 flattering?
\item Which three risks hide in a spread book?
\item How would you backtest the crack fade honestly?
\item When should a location fade stop?
\item Which strategy file needs a fundamental model?
\item How does a spread book relate to chapter~16's systematic macro?
\item In one sentence: what does a spread trader bet on?
\end{enumerate}
\end{problem}

\section{Interview questions}

\begin{interviewq}[$\star$ \iqroles{trader}]\label{iq:s2:commodity-spreads:1}
What is a 3-2-1 crack spread, and who hedges it?
\end{interviewq}

\begin{interviewq}[$\star\star$ \iqroles{researcher}]\label{iq:s2:commodity-spreads:2}
A spread shows strong mean reversion in daily data. How do you check it is not an artefact?
\end{interviewq}

\begin{interviewq}[$\star\star$ \iqroles{trader}]\label{iq:s2:commodity-spreads:3}
Brent-WTI has widened \$10 in a month. Do you fade it?
\end{interviewq}

\begin{interviewq}[$\star\star$ \iqroles{risk}]\label{iq:s2:commodity-spreads:4}
How would you set limits on a book of commodity spreads?
\end{interviewq}

\begin{interviewq}[$\star\star$ \iqroles{developer}]\label{iq:s2:commodity-spreads:5}
How would you store and roll the legs of a three-leg spread in a backtester?
\end{interviewq}

\begin{interviewq}[$\star\star\star$ \iqroles{researcher}]\label{iq:s2:commodity-spreads:6}
Two prices follow the same random walk but are recorded at different times of day, a fraction $f$ of a day apart. Show that the daily changes of
their difference have a first-order autocorrelation of $-1/2$ whatever $f$, and say what happens when the true spread also moves.
\end{interviewq}
